\documentclass[../main.tex]{subfiles}
\begin{document}


\vspace{1em}

Trong những năm gần đây, sự tiến hóa của các phương thức tấn công mạng đã đặt các Trung tâm điều hành an ninh (Security Operations Center - SOC, đơn vị chuyên trách giám sát và bảo vệ hệ thống mạng của tổ chức) vào trạng thái quá tải liên tục. Nhật ký sự kiện hệ thống trên máy trạm (endpoint logs - dữ liệu ghi nhận mọi hoạt động cốt lõi của hệ điều hành), đặc biệt là công cụ Windows System Monitor (Sysmon), đã trở thành nguồn dữ liệu pháp y nền tảng để theo dõi và truy vết các hành vi độc hại sâu bên trong hệ điều hành Windows. Tuy nhiên, sự bùng nổ về khối lượng dữ liệu với hàng triệu sự kiện mỗi ngày trên một máy trạm đã tạo ra một rào cản thông tin khổng lồ, đòi hỏi các giải pháp phân tích tự động thay vì phụ thuộc hoàn toàn vào sức người.

Đồng thời, các tác nhân đe dọa có chủ đích (Advanced Persistent Threat - APT) hiện đại đã thay đổi hoàn toàn chiến thuật để né tránh các hệ thống phòng thủ. Nổi bật nhất là kỹ thuật Living-off-the-Land (LotL) — lạm dụng trực tiếp các tiến trình hệ thống hợp lệ và có sẵn của Windows (như \texttt{powershell.exe}, \texttt{wmic.exe}, \texttt{cmd.exe}) thay vì đưa mã độc lạ vào đĩa cứng. Sự dịch chuyển này làm cho các hệ thống phát hiện dựa trên chữ ký tĩnh truyền thống (như hệ thống quản lý sự kiện SIEM hay giải pháp bảo vệ máy trạm EDR) bị suy giảm hiệu quả nghiêm trọng. Kể cả việc áp dụng các tập luật đối sánh linh hoạt (như bộ quy tắc định dạng mở Sigma \cite{sigmahq2024}) cũng tạo ra tỷ lệ báo động giả (False Positive) cực kỳ cao do không thể phân biệt rạch ròi giữa hành vi quản trị hệ thống hợp lệ và hành vi tấn công ngụy trang. Đứng trước sự giới hạn của đối sánh chữ ký tĩnh, phương pháp Phát hiện bất thường không giám sát (Unsupervised Anomaly Detection - UAD) được kỳ vọng là giải pháp tối ưu nhờ khả năng tự động học hỏi hành vi bình thường và phát hiện các mối đe dọa zero-day mà không cần gán nhãn dữ liệu huấn luyện.

Dù nhận được sự kỳ vọng lớn, thực tế triển khai UAD lại bộc lộ nhiều điểm yếu chí mạng. Nhiều nghiên cứu trong phòng thí nghiệm áp dụng các thuật toán không giám sát truyền thống lên các bộ dữ liệu đơn giản đã báo cáo mức hiệu năng phát hiện rất cao với F1-Score vượt quá 0.85 \cite{ahmed2016survey, alsaheel2021atlas, han2020unicorn}. Tuy nhiên, khi đối mặt với các chiến dịch APT thực tế áp dụng LotL (như APT29), hiệu năng của các mô hình này sụp đổ nghiêm trọng. Một số công trình học thuật gần đây \cite{smiliotopoulos2025detection, ugea_lmd_2025} đã bắt đầu ghi nhận và cảnh báo về giới hạn của học máy không giám sát trước các hành vi ngụy trang tinh vi. Mặc dù vậy, khoảng trống nghiên cứu lớn nhất hiện nay là sự vắng bóng hoàn toàn của một thước đo toán học định lượng để chẩn đoán độ khó nội tại (intrinsic difficulty) của chính tập dữ liệu tấn công. Các đánh giá hiện tại hoàn toàn mang tính kinh nghiệm và phụ thuộc vào thuật toán, dẫn đến việc thiếu cơ sở lý thuyết để giải thích nguyên nhân và thời điểm các mô hình này thất bại. 

Luận văn này tập trung giải quyết bài toán: Làm thế nào để lượng hóa một cách độc lập mức độ hòa trộn hành vi của mã độc ngụy trang và giải thích cơ chế thất bại của các thuật toán không giám sát truyền thống. Nghiên cứu tập trung vào các tiến trình Windows được trích xuất từ Sysmon, biểu diễn trong không gian đặc trưng đa chiều hỗn hợp ($V_{10}$), và đánh giá trên 11 tập dữ liệu (bao gồm BOTSv1 \cite{botsv1}, FIN6 \cite{mitre_fin6}, APT29 \cite{mitre_apt29} và các kịch bản tổng hợp). Quá trình phân tích thực nghiệm được đặt trong giới hạn dữ liệu pháp y tĩnh (offline forensics), qua đó cung cấp một giới hạn trên lý thuyết (theoretical upper bound) về hiệu năng phát hiện trong môi trường lý tưởng. 

Thay vì đề xuất thêm một mô hình học máy hộp đen mới, đồ án mang đến giải pháp định hình lại cách cộng đồng đánh giá bài toán phát hiện mã độc LotL. Nghiên cứu đề xuất khung đo lường Behavioral Blending Score (BBS) với thước đo cốt lõi \texttt{DensPct\_Mahal} — một thước đo độ khó hoàn toàn độc lập với thuật toán, giúp giải thích cơ chế thất bại của các bộ phát hiện truyền thống chính xác hơn các phép đo thông thường. Bên cạnh đó, nghiên cứu khám phá sự tồn tại của các vùng rỗng cục bộ ("Local Pockets") — nơi mã độc lẩn trốn sự giám sát toàn cục, và sử dụng mô hình KDE\_Joint như một bằng chứng khái niệm để khẳng định khả năng vượt qua thiên kiến trung tâm của các thuật toán phân tích cấu trúc.

Để trình bày chi tiết các nội dung trên, ngoài phần Mở đầu và Kết luận, luận văn được cấu trúc thành ba chương chính. Chương 1 tổng hợp cơ sở lý thuyết về Windows Event Log, kỹ thuật tấn công APT, các khái niệm bất thường và phương pháp phát hiện không giám sát. Chương 2 đi sâu vào phân tích toán học, xây dựng khung đo lường BBS và thiết kế môi trường thử nghiệm. Cuối cùng, Chương 3 trình bày các kịch bản đánh giá thực nghiệm, đối chiếu kết quả điểm số với hiệu năng thực tế của các mô hình UAD, từ đó đưa ra các phân tích và giải thích cặn kẽ về giới hạn của AI trước các cuộc tấn công ngụy trang.

\end{document}
